\section{Falsifiability}

This section addresses the falsifiability of the claims made about STOP as a diagnostic outcome. Its purpose is to specify the conditions under which the statements advanced in this paper could be shown to be false, thereby establishing their scientific admissibility. No operational guidance or prescriptive criteria are introduced.

The central claim of the paper is not that STOP occurs in particular enterprises, but that STOP can constitute a formally valid diagnostic outcome under specified modeling conditions. Accordingly, falsifiability is assessed at the level of the model and its logical consequences, rather than at the level of individual organizational decisions. This approach is consistent with inquiry-oriented views of systems analysis, in which refutation applies to representations and their implications rather than to observed actions \cite{churchman1971}.

The framework would be falsified if it could be demonstrated that, within a well-specified enterprise model, a STOP outcome is declared while an admissible continuation state nonetheless exists under the same model and constraints. Such a case would show that STOP, as defined, does not correctly track admissibility and would invalidate the proposed interpretation. This falsification criterion relies on the constraint-based notion of admissible state spaces \cite{ashby1956}.

Conversely, the framework would also be falsified if continuation were shown to be analytically unavoidable in all admissible enterprise models, regardless of their constraints or structure. Such a result would contradict the core claim that the admissible continuation space may, under certain conditions, be empty.

A further falsification condition concerns model dependence. If STOP outcomes were shown to be invariant under arbitrary changes of the enterprise model, this would contradict the proposition that STOP is relative to a specific modeling frame. Model-invariant STOP would undermine the claim that STOP reflects admissibility rather than an absolute property of the enterprise.

Importantly, falsifiability does not require empirical observation of STOP decisions in practice. It requires only the possibility of constructing or identifying a model instance that contradicts the stated properties of STOP as a diagnostic outcome. The relevant tests are therefore logical and structural rather than empirical in a narrow sense, consistent with established accounts of modeling and anticipatory systems \cite{rosen1985}.

By specifying these falsification conditions, the framework positions STOP as a scientifically testable concept. The claims advanced are not insulated from refutation; they stand or fall with the coherence of the relationship between model, admissibility, and diagnostic outcome.
